\chapter{Trabalhos Relacionados}
\label{cap:trabalhos-relacionados}

\section{Abordagens baseadas em busca}

Antes de redes neurais mais profundas serem amplamente utilizadas para produção de projetos envolvendo o \textit{Texas Hold'Em}, eram utilizadas abordagens mais voltadas para busca em árvore, como feito por \citeonline{kleij2010thesis} e por \citeonline{vdBroeck2009mcts}, que utilizaram o algoritmo de \textit{Monte Carlo} para explorar os estados do jogo, porém como ambos focavam majoritariamente na modelagem e busca em árvore de maneira mais bruta, enquanto este trabalho é mais focado na tomada de decisões em si, porém o \textit{Monte Carlo} é usado de forma semelhante nas fase de \textit{Pré-Flop} e de \textit{Pós-Flop} para calcular o valor de equidade.

\section{Avanço em \textit{Deep Reinforcement Learning}}

Com o passar do tempo, técnicas de Aprendizado por Reforço Profundo se tornaram mais utilizadas para desenvolvimento de agentes inteligentes capazes de jogar o \textit{Texas Hold'Em}, como feito por \citeonline{moravvcik2017deepstack}, que desenvolveu um agente capaz de vencer jogadores profissionais de poquêr \textit{Heads-Up}, outra categoria do jogo quem consiste em um contra um direto, utilizando clusters poderosos carregados de dados, enquanto este trabalho é mais focado na intuição dos jogos.

Enquanto a maioria dos trabalhos realizados para criação de agentes para jogar poquêr jogam apenas a categoria \textit{Heads-Up}, e enfrentam inúmeros problemas de ótimo local, assim como \citeonline{mnih2015human} tentou desenvolver um agente utilizando apenas uma \gls{DQN} puro para jogar o \textit{Texas Hold'Em}, mas atingiu ótimos locais de desistências infinitas assim como este trabalho, tendo isso em vista \citeonline{heinrich2016deep} desenvolveu um agente capaz de jogar \textit{Texas Hold'Em} utilizando táticas para contornar as fraquezas do \gls{DQN}, enquanto a tática utilizada neste trabalho foi a de penalidade por jogar passivamente.

Também pode ser citado o artigo de \citeonline{brown2019superhuman} que desenvolveu um agente que se sobressaiu em uma mesa com 6 jogadores, demonstrando um domínio completo sobre ambientes multi-jogadores, o agente jogava de maneira que não permitia que os outros jogadores explorassem suas falhas, tendo um comportamento similar ao agente desenvolvido neste trabalho, que explora a falha dos outros jogadores para tomar as melhores decisões possíveis.

\section{Diferença arquitetural}

Enquanto os trabalhos citados anteriormente dispunham de um extenso maquinário para o desenvolvimento dos agentes, este trabalho se voltou mais para o fator de exploração do agente, fazendo com que o próprio aprendesse com suas próprias ações, além de utilizar modulagem de recompensa para deixar o agente mais preciso, o penalizando por fazer uma jogada desnecessariamente passiva, ou o recompensando por fazer uma jogada embasadamente agressiva, fazendo com que o jogador tenda a seguir mais fielmente a dinâmica do \textit{Texas Hold'Em}.

Também pode ser citada a questão de desempate, que costuma ser resolvida por buscas em árvores complexas e mais custosas, este projeto desenvolveu um motor de base 15 capaz de efetuar os desempates de forma concisa e coerente, deixando o programa relativamente mais rápido. A linguagem utilizada para o desenvolvimento também é diferente da maioria utilizada, já que foi utilizado o \textit{JavaScript}, utilizando do \textit{TensorFlow.js} e \textit{Node.js} como uma alternativa no lugar da comumente usada \textit{Python}.

